% ============================================================
% EJERCICIO 3 — Informe de laboratorio completo (Módulos 1 y 2)
% Tema: «Circuito Amplificador Operacional Inversor»
% Tareas: 1) pida a una IA la estructura (guarde el prompt);
%         2) complete los TODO; 3) compile 2 veces.
% Datos de la tabla: simulados o medidos por usted (indique cuál).
% PROMPT UTILIZADO: TODO
% ============================================================
\documentclass[12pt,a4paper]{article}
\usepackage[utf8]{inputenc}
\usepackage[T1]{fontenc}
\usepackage[spanish]{babel}
\usepackage{amsmath,amssymb}
\usepackage{booktabs}
\usepackage{siunitx}
\sisetup{locale=FR, output-decimal-marker={,}}
\usepackage{circuitikz}
\usepackage[hidelinks]{hyperref}
\usepackage{float}

\begin{document}
% TODO Portada (titlepage): universidad, facultad, título, autor, fecha
\tableofcontents
\newpage

\section{Objetivos}
% TODO lista de 2 objetivos

\section{Marco teórico}
% TODO al menos 2 ecuaciones numeradas: ganancia -Rf/R1 y saturación

\section{Circuito}
% TODO figura con circuitikz (op amp), caption y label; cítela con \ref

\section{Resultados}
% TODO tabla Vin vs Vout (5 filas) con booktabs + siunitx; cítela con \ref

\section{Conclusiones}
% TODO

% TODO Bibliografía: 2 fuentes REALES que usted haya abierto
\end{document}
% RÚBRICA (20 pts): portada+índice (3) · ecuaciones (4) · figura circuitikz (4)
% · tabla correcta (3) · referencias cruzadas (2) · bibliografía verificada (2)
% · prompt y revisión crítica de la IA (2)
